%%%%%%%%%%%%%%%%%
% CV tailored for Thales - Responsable Offres Logicielles F/H
% Form inspired by the user's Intheair CV structure: clear context/objective blocks,
% concise bullets, strong role positioning.
%%%%%%%%%%%%%%%%%

\documentclass[8pt,a4paper,ragged2e,withhyper]{altacv}

\geometry{left=1.25cm,right=1.25cm,top=.5cm,bottom=1.cm,columnsep=1.2cm}

\usepackage{paracol}
\usepackage{fontawesome5}
\usepackage{hyperref}

\iftutex 
  \setmainfont{Roboto Slab}
  \setsansfont{Lato}
  \renewcommand{\familydefault}{\sfdefault}
\else
  \usepackage[rm]{roboto}
  \usepackage[defaultsans]{lato}
  \renewcommand{\familydefault}{\sfdefault}
\fi

\definecolor{SlateGrey}{HTML}{2E2E2E}
\definecolor{LightGrey}{HTML}{666666}
\definecolor{DarkPastelRed}{HTML}{8AC0D9}
\definecolor{PastelRed}{HTML}{00B0DA}
\definecolor{GoldenEarth}{HTML}{B4AD0C}

\colorlet{name}{black}
\colorlet{tagline}{PastelRed}
\colorlet{heading}{DarkPastelRed}
\colorlet{headingrule}{GoldenEarth}
\colorlet{subheading}{PastelRed}
\colorlet{accent}{PastelRed}
\colorlet{emphasis}{SlateGrey}
\colorlet{body}{LightGrey}

\definecolor{violet}{HTML}{5A2582}
\definecolor{rouge}{HTML}{F53B3B}
\definecolor{noir}{HTML}{00232C}
\definecolor{bleu}{HTML}{00B0DA}
\definecolor{rose}{HTML}{F831A0}
\definecolor{jaune}{HTML}{FFEC00}
\definecolor{vert}{HTML}{34AD6C}

\renewcommand{\namefont}{\Huge\rmfamily\bfseries}
\renewcommand{\personalinfofont}{\footnotesize}
\renewcommand{\cvsectionfont}{\LARGE\rmfamily\bfseries}
\renewcommand{\cvsubsectionfont}{\large\bfseries}
\renewcommand{\cvItemMarker}{{\small\textbullet}}
\renewcommand{\cvRatingMarker}{\faCircle}

\input{../../_shared/pubs-num.tex}
\addbibresource{../../_shared/sample.bib}

\begin{document}

\name{BOILEAU Guillaume}
\tagline{Responsable offres logicielles / Coordination technique | Devis, systèmes complexes, IVVQ \& reporting}

\photoL{2.8cm}{../../_shared/moi}

\personalinfo{%
  \email{guillaume.boileaupro@gmail.com}
  \phone{+33 6 59 94 85 84}
  \location{Alpes-Maritimes, FRANCE}
  \linkedin{guillaume-boileau-891b81265}
  \researchgate{Guillaume-Boileau-2}
  \orcid{0000-0002-3576-6968}
}

\makecvheader
\vspace{-0.3cm}

\AtBeginEnvironment{itemize}{\small}
\columnratio{0.64}

\begin{paracol}{2}

\cvsection{Experience}

\cvevent{\textbf{Ingénieur intégration logiciel -- Interfaces, support opérationnel \& coordination agile}}{VEV Platform Services France}{2025 -- 2026}{Valbonne, France}

\textbf{Contexte :} Environnement industriel logiciel pour une plateforme CPMS liée à des infrastructures de recharge de véhicules électriques, avec services backend, interfaces applicatives, données opérationnelles et équipements terrain.

\textbf{Objectif :} Contribuer au développement, à l’intégration et au suivi technique de services logiciels connectés, en assurant la cohérence entre flux de données, comportements applicatifs, contraintes opérationnelles et priorités de livraison.

\begin{itemize}
  \item \textbf{Livraison logicielle :} Contribution à des composants TypeScript, Angular et Node.js intégrés dans une plateforme opérationnelle.
  \item \textbf{Intégration système :} Travail sur les interfaces entre services logiciels, données opérationnelles, équipements connectés et contraintes de production.
  \item \textbf{Suivi opérationnel :} Analyse de flux FTP, interruptions d’échanges, incohérences de données et comportements applicatifs inattendus.
  \item \textbf{Validation technique :} Vérification de la cohérence entre données applicatives, comportement terrain et attendus fonctionnels.
  \item \textbf{Support production :} Investigation d’incidents, diagnostic technique, formalisation des constats et contribution aux actions de fiabilisation.
  \item \textbf{Coordination agile :} Suivi des tâches, priorisation, backlog, points d’avancement et coordination via Zenhub.
  \item \textbf{Reporting technique :} Synthèse d’observations, d’anomalies et de points bloquants pour faciliter les arbitrages.
\end{itemize}

\divider

\cvevent{\textbf{Ingénieur R\&D senior -- Coordination technique, logiciel scientifique \& validation}}{Laboratoire Lagrange, Observatoire de la Côte d’Azur, CNES}{2023 -- 2025}{Nice, France}

\textbf{Contexte :} Activités d’ingénierie et de R\&D pour la mission spatiale LISA ESA/NASA, dans un environnement international impliquant simulation, développement Python, traitement de données, intégration de modèles et validation technique.

\textbf{Objectif :} Structurer, intégrer, comparer et valider des contributions logicielles et scientifiques complexes afin de produire des résultats exploitables, traçables et présentables en revue collaborative.

\begin{itemize}
  \item \textbf{Coordination technique :} Interface avec des contributeurs scientifiques, logiciels et institutionnels dans un environnement spatial international.
  \item \textbf{Lots logiciels et livrables :} Structuration de workflows Python, versions, configurations, jeux de données, sorties de simulation et résultats d’analyse.
  \item \textbf{Jugement technique :} Analyse des contributions, identification des points essentiels, consolidation des résultats et formulation de synthèses exploitables.
  \item \textbf{Validation / IVVQ :} Définition de contrôles de cohérence, règles de comparaison, méthodes de vérification et critères d’acceptabilité des résultats.
  \item \textbf{Pilotage de jalons techniques :} Préparation d’éléments de revue, consolidation des contributions et suivi des versions pour sécuriser les décisions.
  \item \textbf{Risques \& cohérence technique :} Identification des écarts entre modèles, analyse d’impact, formalisation des points bloquants et proposition d’actions correctives.
  \item \textbf{Documentation et capitalisation :} Utilisation de Confluence pour organiser la documentation, la traçabilité, le suivi de validation et le partage de connaissances.
  \item \textbf{Plateforme logicielle :} Développement de CATpyLISA, plateforme Python pour l’intégration, la comparaison, la validation et la revue technique de modèles.
  \textcolor{DarkPastelRed}{\href{https://gitlab.oca.eu/gboileau/lisa_l3_tools}{\bf GitLab: CATpyLISA}}
\end{itemize}

\divider

\cvevent{\textbf{Data Scientist -- Modélisation, bruit \& analyse de performance}}{Universiteit Antwerpen, Belgium}{2021 -- 2023}{Antwerp, Belgium}

\textbf{Contexte :} Études de performance pour instruments de précision et détecteurs de nouvelle génération, combinant analyse de données, bruit environnemental, modélisation statistique et validation de résultats.

\textbf{Objectif :} Développer des workflows robustes et reproductibles pour identifier, caractériser et réduire des sources affectant la qualité de mesure et la sensibilité de systèmes complexes.

\begin{itemize}
  \item \textbf{Analyse de performance :} Évaluation de la qualité de mesure, des limites instrumentales et de l’impact de différentes sources de bruit.
  \item \textbf{Traitement du signal :} Identification et caractérisation de contributions de bruit dans des mesures complexes.
  \item \textbf{Modélisation statistique :} Développement de méthodes d’analyse et d’inférence bayésienne, notamment par chaînes MCMC.
  \item \textbf{Configurations système :} Études de différentes configurations instrumentales et de différents niveaux de bruit pour la détection de signaux très ténus.
  \item \textbf{Canaux témoins \& ML :} Structuration d’approches de réduction de bruits non stationnaires avec canaux témoins, machine learning et deep learning.
  \item \textbf{Validation de résultats :} Analyse de sensibilité, vérification de cohérence, figures de mérite et évaluation des limites de performance.
  \item \textbf{Reporting technique :} Production de résultats traçables, de synthèses techniques et de supports pour parties prenantes internationales.
\end{itemize}

\divider

\cvevent{\textbf{Ingénieur R\&D junior -- Signal, simulation \& validation logicielle}}{ARTEMIS, Observatoire de la Côte d’Azur}{2018 -- 2021}{Nice, France}

\textbf{Contexte :} Activités de R\&D liées à l’analyse de signaux faibles, à la simulation numérique, aux diagnostics de bruit et à la préparation de missions spatiales et d’instruments de détection.

\textbf{Objectif :} Développer des méthodes d’analyse, automatiser des simulations et produire des résultats exploitables dans un environnement scientifique et technique exigeant.

\begin{itemize}
  \item \textbf{Simulation numérique :} Développement d’approches de simulation pour environnements de signaux complexes et jeux de données de grande taille.
  \item \textbf{Analyse de données :} Mise en place de méthodes pour analyser et séparer des signaux faibles et superposés.
  \item \textbf{Traitement du signal :} Études de caractérisation, de décomposition et de comparaison de signaux dans des contextes bruités.
  \item \textbf{Validation technique :} Vérification de cohérence, comparaison de résultats et études de sensibilité.
  \item \textbf{Investigation technique :} Analyse de sources de bruit et diagnostics par corrélations entre capteurs.
  \item \textbf{Documentation :} Formalisation des méthodes, hypothèses, résultats et conclusions pour revue collaborative.
\end{itemize}

\switchcolumn

\cvsection{Profile}

\textbf{Ingénieur docteur orienté offres logicielles, coordination technique et systèmes complexes, avec une expérience solide en développement Python, intégration logicielle, validation, analyse de performance, reporting et documentation technique en environnement spatial international.}

Mon parcours combine livraison logicielle, structuration de contributions techniques, synthèse d’informations complexes, validation de résultats, coordination transverse, capitalisation documentaire et support à la décision. J’ai travaillé dans des environnements exigeants associant logiciel, données, simulation, signal, instrumentation et contraintes de systèmes critiques.

\begin{itemize}
  \item Capacité à piloter et consolider des contributions logicielles dans des systèmes complexes.
  \item Expérience en synthèse technique, arbitrage, structuration de livrables, reporting et préparation de revues.
  \item Culture forte de la validation, de l’IVVQ, de la traçabilité, des exigences et de la qualité des résultats.
  \item Profil rigoureux, engagé, à l’aise dans les interfaces entre équipes logiciel, système, métier et management.
\end{itemize}

\cvsection{Languages}

\cvskill{English}{4}
\divider
\cvskill{Spanish}{2}
\divider

\medskip

\cvsection{Education}

\cvevent{PhD in Physics}{Université Côte d’Azur}{2018 -- 2021}{Nice, France}

\divider

\cvevent{Selective Graduate Program in Fundamental Physics (Magistère)}{Université Paris-Saclay}{2016 -- 2018}{Orsay, France}

\divider

\cvevent{MSc in Astronomy and Astrophysics}{Observatoire de Paris / Université Paris-Saclay}{2017 -- 2018}{Paris, France}

\divider

\cvevent{BSc in Fundamental Physics}{Université Paris-Saclay}{2015 -- 2016}{Orsay, France}

\cvsection{Professional Training}

\cvevent{Machine Learning in Python with scikit-learn}{FUN MOOC -- Inria}{2026}{France Université Numérique}

\small{
Predictive modelling, supervised learning, model evaluation, cross-validation, metrics, pipelines and practical machine learning workflows with Python and scikit-learn.
}

\divider

\cvevent{Supervised Machine Learning: Regression \& Classification}{DeepLearning.AI -- Andrew Ng}{2020}{Coursera}

\divider

\cvevent{Recherche reproductible : principes méthodologiques pour une science transparente}{FUN MOOC -- Inria}{2025}{France Université Numérique}

\divider

\cvevent{Continuous Professional Training -- Software, Data, AI and Systems}{OpenClassrooms}{2024 -- 2026}{}

\small{
Python, SQL, Machine Learning, AI, Linux, JavaScript, TypeScript, Angular, Node.js, Django, REST APIs, C/C++, web development, algorithms and software engineering fundamentals.
}

\end{paracol}

\cvsection{Projects}

\cvevent{CATpyLISA -- Plateforme Python d’intégration, comparaison et validation}{Mission spatiale LISA ESA/NASA}{2023 -- 2025}{}

Développement d’une plateforme Python dédiée à l’intégration de modèles, à la structuration de données, à la comparaison de résultats, à la validation technique et à la revue de workflows liés à la mission LISA.

\textbf{Rôle :} développement Python, conception de pipelines, structuration de données, logique de validation, contrôles de cohérence, documentation technique, consolidation de contributions et coordination avec plusieurs contributeurs.

\textbf{Compétences mobilisées :} Python, intégration de modèles, validation, traitement du signal, analyse de performance, simulation, documentation, GitLab, Confluence, environnement spatial international.

\divider

\cvevent{VEV -- Livraison logicielle, interfaces \& support opérationnel}{Industrial software / Connected infrastructure}{2025 -- 2026}{}

Contribution au développement, à l’intégration, au suivi opérationnel et à la validation de services logiciels connectés à des infrastructures terrain.

\textbf{Rôle :} développement logiciel, investigation d’incidents, analyse de flux de données, support opérationnel, suivi de tickets, coordination agile et reporting technique.

\textbf{Compétences mobilisées :} TypeScript, Angular, Node.js, FTP, Linux, Git, Zenhub, analyse d’anomalies, support production.

\divider

\cvevent{Workflows de simulation, performance et analyse de bruit}{Instrumentation scientifique et préparation de missions spatiales}{2018 -- 2025}{}

Développement de workflows numériques et statistiques pour analyser des signaux faibles, caractériser des contributions de bruit, comparer des résultats et évaluer la robustesse de modèles dans des environnements complexes.

\textbf{Rôle :} analyse de performance, simulation, traitement du signal, caractérisation de bruit, études de sensibilité, diagnostics techniques, reporting et support à la décision.

\textbf{Compétences mobilisées :} traitement du signal, simulation numérique, analyse spectrale, Python, C, validation, analyse de sensibilité, documentation technique.

\cvsection{Compétences clés}

\vspace{-.3cm}

\cvsubsection{Offres logicielles, devis \& pilotage}

{\small
\cvtag{Pilotage d’offres logicielles}
\cvtag{Coordination de réponse}
\cvtag{Consolidation technique}
\cvtag{Devis logiciel -- ramp-up}
\cvtag{Chiffrage -- ramp-up}
\cvtag{Jugement technique}
\cvtag{Synthèse décisionnelle}
\cvtag{Reporting d’avancement}
\cvtag{Dashboard}
\cvtag{KPIs}
\cvtag{Capitalisation}
\cvtag{REX}
\cvtag{Suivi financier -- awareness}
\cvtag{Rigueur}
}

\divider\smallskip
\vspace{-.3cm}

\cvsubsection{Coordination ingénierie \& systèmes complexes}

{\small
\cvtag{Coordination technique}
\cvtag{Pilotage transverse}
\cvtag{Interfaces métiers}
\cvtag{Interfaces système / logiciel}
\cvtag{Responsabilité de livrables}
\cvtag{Suivi de jalons}
\cvtag{Gestion de priorités}
\cvtag{Préparation de revues}
\cvtag{Support décisionnel}
\cvtag{Leadership technique}
\cvtag{Communication transverse}
\cvtag{Anglais technique}
}

\divider\smallskip
\vspace{-.3cm}

\cvsubsection{Logiciel, intégration \& livraison}

{\small
\cvtag{Logiciels intégrés}
\cvtag{Systèmes complexes}
\cvtag{Intégration logicielle}
\cvtag{Développement Python}
\cvtag{TypeScript}
\cvtag{Angular}
\cvtag{Node.js}
\cvtag{REST API}
\cvtag{SQL}
\cvtag{Bash}
\cvtag{Linux}
\cvtag{Git}
\cvtag{GitLab}
\cvtag{Docker}
\cvtag{CI/CD}
\cvtag{FTP}
}

\divider\smallskip
\vspace{-.3cm}

\cvsubsection{IVVQ, qualité \& validation}

{\small
\cvtag{IVVQ logiciel}
\cvtag{IVVQ système}
\cvtag{Vérification / Validation}
\cvtag{Contrôles de cohérence}
\cvtag{Tests fonctionnels}
\cvtag{Non-régression}
\cvtag{Analyse d’anomalies}
\cvtag{Root cause analysis}
\cvtag{Diagnostic technique}
\cvtag{Traçabilité}
\cvtag{Documentation qualité}
\cvtag{Confluence}
\cvtag{Jira}
\cvtag{Zenhub}
}

\divider\smallskip
\vspace{-.3cm}

\cvsubsection{Signal, simulation \& domaines techniques}

{\small
\cvtag{Traitement du signal}
\cvtag{Acoustique -- ramp-up}
\cvtag{Sonar -- ramp-up}
\cvtag{Simulation numérique}
\cvtag{Analyse de performance}
\cvtag{Analyse de bruit}
\cvtag{Signaux faibles}
\cvtag{Low-SNR}
\cvtag{Modélisation}
\cvtag{Data quality}
\cvtag{Machine Learning}
\cvtag{Systèmes critiques}
\cvtag{Spatial}
\cvtag{Défense -- habilitation compatible}
}

\divider\smallskip
\vspace{-.3cm}

\cvsubsection{Outils \& monitoring}

{\small
\cvtag{Python}
\cvtag{Pandas}
\cvtag{NumPy}
\cvtag{SciPy}
\cvtag{Matplotlib}
\cvtag{scikit-learn}
\cvtag{Jupyter}
\cvtag{Power BI}
\cvtag{OpenSearch}
\cvtag{Sentry}
\cvtag{Confluence}
\cvtag{Jira}
\cvtag{Zenhub}
}

\end{document}
